\section{Research challenges}
\label{sec:Challenges}

Although hybrid cloud provides flexibility and scalability, it also introduces some significant challenges in managing secure delivery and infrastructure operations. In addition, the hybrid cloud requires organizations to reconcile software deployment, security policies, and compliance while maintaining system performance. These challenges are identified as follows.

\vspace{0.3cm}

\textit{Challenge 1: Managing heterogeneous infrastructure}.
Hybrid cloud refers to multiple diverse computing platforms, including on-premises infrastructure, private cloud, and public cloud, where each one has different architectures, interfaces, and security models. Consequently, the transition from the traditional cloud model to the hybrid cloud model requires maintaining consistent operational procedures and configuration policies, which is considered to be complex and time-consuming. \cite{Liu2011NISTRA} also states that the hybrid cloud needs interoperability and portability among many distinct cloud infrastructures while preserving their individual characteristics. As a result, operational teams must manually manage diverse infrastructures that often consist of different methods, mechanisms, and security policies.

Moreover, \cite{Flexera2025StateCloud}, \cite{HashiCorp2025StateCloud} agree that as organizations expand their models, the complexity of models escalates, increasing administrative overhead and the likelihood of security misconfigurations. This, by far, makes consistency across heterogeneous infrastructure increasingly difficult. \cite{HashiCorp2025StateCloud} highlights that the increasing complexity in cloud infrastructure has become one of the primary concerns for most organizations. In addition, differences in configuration policies, identity management, and resource allocation increase administrative overhead and configuration inconsistencies that may expose security vulnerabilities.

\vspace{0.3cm}

\textit{Challenge 2: Integrating continuous security into the software lifecycle}.
Traditional security operations are usually performed during the later stages of software development, making vulnerabilities more expensive and difficult to detect and remediate. Although this approach is capable of maintaining security standards for conventional software development models, it may expose serious vulnerabilities in modern cloud-native applications, which are continuously developed and deployed through automated CI/CD pipelines. Thus, \cite{NISTSSDF2022} emphasizes that security operations should be continuously incorporated into software development and deployment processes to reduce software vulnerabilities and improve software integrity.

In hybrid cloud infrastructure, applications can be deployed across multiple infrastructures with different security requirements. Therefore, Security operations must perform not only on application source code but also on software dependencies, configuration policies, and resource allocation. Performing these assessments independently or only before production deployment increases the risk that vulnerabilities remain undetected until later stages, where remediation is more expensive and difficult.

Furthermore, organizations must integrate automated security operations into existing development workflows without interfering with development productivity. Additionally, Continuous vulnerability scanning and policy validation should operate alongside continuous integration and deployment processes while maintaining acceptable execution time. Consequently, many organizations try to balance between rapid software delivery and continuous security in hybrid cloud infrastructure \cite{NISTSSDF2022}.

\vspace{0.3cm}

\textit{Challenge 3: Assessing security findings and prioritizing risks}.
Modern software applications are usually developed with multiple security evaluation tools. These include static application security testing, software composition analysis, dependency vulnerability scanners, and infrastructure configuration analysis tools. Each tool focuses on detecting different categories of vulnerabilities and produces systematic reports with its own severity metrics and confidence levels. Although these tools cover each security category, they also generate diverse security findings that are difficult to evaluate consistently and prioritize risks.

In practice, security engineers manually analyze findings from multiple scanners, eliminate duplicate reports, and compare vulnerabilities to decide multiple actions for each risk. However, the manual process becomes more challenging as project size grows, which can slow down the whole deployment. In hybrid cloud infrastructures, where applications are updated and deployed frequently on multiple clouds, delayed or inconsistent prioritization may leave critical vulnerabilities unresolved while less significant issues receive unnecessary attention.

\cite{NISTSSDF2022} recommends integrating security operations into a software development process to support continuous vulnerability management and prioritization. Nevertheless, effectively combining many different security findings with distinct severity metrics into meaningful risk assessments remains a significant challenge. Thus, organizations require systematic approaches capable of combining multiple security reports and prioritizing vulnerabilities according to their overall operational risk.

\vspace{0.3cm}

\textit{Challenge 4:Establishing an integrated SysSecOps process}.
The complexity of hybrid cloud environments requires close contact between infrastructure management, software delivery, and security operations. However, these activities are performed by separate teams using independent tools and workflows. Specifically, infrastructure administrators concentrate primarily on provisioning and maintaining computing resources, development teams emphasize rapid software delivery, while security teams can only perform security assessments independently after implementation. Consequently, these different processes create communication gaps, which delay security reports and reduce the overall efficiency of software delivery processes.

Although DevOps has significantly reduced the communication gaps between development and security teams, security assessments remain inconsistent across many organizations. Moreover, security operations are usually known as isolated scanning stages rather than as continuous processes from infrastructure management, software development, deployment, to maintenance. As a consequence, when hybrid cloud infrastructures continue to grow in scale and complexity, security operations become increasingly difficult to manage while maintaining consistent security policies and regulatory compliance \cite{HashiCorp2025StateCloud}, \cite{Liu2011NISTRA}.